\subsection{Лемма о разности элементов полукольца}

\begin{lemma}
    Если $A, B_1, B_2, \dots, B_m \in S$, то существуют такие $A_1, A_2, \dots, A_n \in S$, что:
    \[
        A \setminus B_1 \setminus B_2 \setminus \dots \setminus B_m = \bigsqcup_{i=1}^{n} A_i.
    \]
\end{lemma}

\begin{Proof}[Доказательство леммы]
    Доказывать будем по индукции по $m$. %[cite: 9]
    
    \textbf{База индукции ($m=1$):} \\
    Утверждение $A \setminus B_1 = \bigsqcup_{i=1}^{n} A_i$ выполняется по определению полукольца (свойство разности). %[cite: 9]

    \textbf{Шаг индукции:} \\
    Пусть утверждение верно для $m$. Докажем для $m+1$. %[cite: 9]
    \[
        A \setminus B_1 \setminus \dots \setminus B_m \setminus B_{m+1} = \left( \bigsqcup_{i=1}^{n} A_i \right) \setminus B_{m+1} = \bigsqcup_{i=1}^{n} (A_i \setminus B_{m+1}).
    \]
    Так как $A_i \in S$ и $B_{m+1} \in S$, по определению полукольца разность $A_i \setminus B_{m+1}$ представима в виде $\bigsqcup_{j=1}^{k_i} C_{ij}$, где $C_{ij} \in S$. %[cite: 9]
    Тогда:
    \[
        \bigsqcup_{i=1}^{n} (A_i \setminus B_{m+1}) = \bigsqcup_{i=1}^{n} \left( \bigsqcup_{j=1}^{k_i} C_{ij} \right) = \bigsqcup_{i, j} C_{ij}.
    \]
    Получили дизъюнктивное объединение конечного числа элементов из $S$. Лемма доказана. %[cite: 9]
\end{Proof}

\begin{Proof}[{{Продолжение доказательства теоремы (проверка, что $K$ --- кольцо)}}]
    \begin{enumerate}[label=\arabic*)]
        \item $\varnothing \in K$, так как $S \subset K$ и $\varnothing \in S$. %[cite: 10]
        \item Пусть $P, Q \in K$. Докажем, что $P \cap Q \in K$ и $P \triangle Q \in K$. %[cite: 10]
        
        По определению $K$, элементы $P$ и $Q$ имеют вид:
        \[
            P = \bigsqcup_{i=1}^{n} P_i, \quad Q = \bigsqcup_{j=1}^{m} Q_j, \qquad P_i, Q_j \in S.
        \]
        
        \textbf{Пересечение:}
        \[
            P \cap Q = \left( \bigsqcup_{i=1}^{n} P_i \right) \cap \left( \bigsqcup_{j=1}^{m} Q_j \right) = \bigsqcup_{i=1}^{n} \bigsqcup_{j=1}^{m} (P_i \cap Q_j).
        \]
        Обозначим $D_{ij} = P_i \cap Q_j$. Так как $S$ замкнуто относительно пересечений, $D_{ij} \in S$. Значит, $P \cap Q = \bigsqcup_{i,j} D_{ij} \in K$. %[cite: 10]

        \textbf{Симметрическая разность:} \\
        Рассмотрим сначала разность $P \setminus Q$: %[cite: 11]
        \[
            P \setminus Q = \left( \bigsqcup_{i=1}^{n} P_i \right) \setminus \left( \bigsqcup_{j=1}^{m} Q_j \right) = \bigsqcup_{i=1}^{n} (P_i \setminus Q_1 \setminus Q_2 \setminus \dots \setminus Q_m).
        \]
        По доказанной лемме, каждый элемент $P_i \setminus Q_1 \setminus \dots \setminus Q_m$ представим как $\bigsqcup_{\ell=1}^{W_i} F_{i\ell}$, где $F_{i\ell} \in S$. %[cite: 11]
        Отсюда $P \setminus Q = \bigsqcup_{i, \ell} F_{i\ell} \in K$.
        
        Аналогично расписывается $Q \setminus P$:
        \[
            Q \setminus P = \bigsqcup_{j=1}^{m} (Q_j \setminus P_1 \setminus \dots \setminus P_n) = \bigsqcup_{j=1}^{m} \bigsqcup_{k=1}^{T_j} E_{jk} \in K, \quad E_{jk} \in S. %[cite: 11]
        \]
        
        Тогда симметрическая разность:
        \[
            P \triangle Q = (P \setminus Q) \sqcup (Q \setminus P) = \left( \bigsqcup_{i, \ell} F_{i\ell} \right) \sqcup \left( \bigsqcup_{j, k} E_{jk} \right) \in K. %[cite: 11]
        \]
    \end{enumerate}
    Мы доказали, что $K$ --- кольцо. Следовательно, $R(S) = K$. Теорема полностью доказана. %[cite: 11]
\end{Proof}

\subsection{$\sigma$-кольца и $\delta$-кольца}

\begin{definition}[$\delta$-кольцо]
    $\delta$-кольцом называется кольцо, замкнутое относительно счетного пересечения:
    \[
        A_1, A_2, \dots \in R \implies \bigcap_{i=1}^{\infty} A_i \in R.
    \] %[cite: 12]
\end{definition}

\begin{definition}[$\sigma$-кольцо]
    $\sigma$-кольцом называется кольцо, замкнутое относительно счетного объединения:
    \[
        A_1, A_2, \dots \in R \implies \bigcup_{i=1}^{\infty} A_i \in R.
    \] %[cite: 12]
\end{definition}

\begin{exercise}[Свойство вложенности]
    Любое $\sigma$-кольцо является $\delta$-кольцом ($\sigma\text{-кольцо} \implies \delta\text{-кольцо}$). %[cite: 12]
\end{exercise}

\begin{Proof}
    Пусть $R$ --- $\sigma$-кольцо и $A_1, A_2, \dots \in R$. Рассмотрим их счетное пересечение. Воспользуемся законами де Моргана: %[cite: 12]
    \[
        \bigcap_{i=1}^{\infty} A_i = \bigcap_{i=1}^{\infty} (A_1 \cap A_i) = A_1 \setminus \left( A_1 \setminus \bigcap_{i=2}^{\infty} A_i \right) = A_1 \setminus \bigcup_{i=2}^{\infty} (A_1 \setminus A_i).
    \]
    Проанализируем правую часть: %[cite: 12]
    \begin{itemize}
        \item $A_1 \setminus A_i \in R$, так как $R$ --- кольцо (замкнуто относительно разности).
        \item $\bigcup_{i=2}^{\infty} (A_1 \setminus A_i) \in R$, так как $R$ --- $\sigma$-кольцо (замкнуто относительно счетного объединения).
        \item Итоговая разность $A_1 \setminus \left( \dots \right)$ также принадлежит $R$, так как кольцо замкнуто относительно разности.
    \end{itemize}
    Следовательно, $\bigcap_{i=1}^{\infty} A_i \in R$, что и требовалось доказать. %[cite: 12]
\end{Proof}

\begin{example}
    Обратное утверждение неверно: $\delta$-кольцо $\not\implies$ $\sigma$-кольцо. %[cite: 12]
    \\
    \\
    В качестве контрпримера рассмотрим систему всех ограниченных подмножеств $\R$. %[cite: 12] 
    Эта система образует $\delta$-кольцо, поскольку пересечение любого (в том числе счетного) числа ограниченных множеств остается ограниченным. Однако она не является $\sigma$-кольцом: счетное объединение ограниченных множеств может дать неограниченное множество. 
\end{example}